\documentclass[12pt,a4paper]{article}
\usepackage[utf8]{inputenc}
\usepackage[T1]{fontenc}
\usepackage{amsmath,amssymb,amsthm}
\usepackage{geometry}
\usepackage{hyperref}
\usepackage{booktabs}

\geometry{margin=1in}

\newtheorem{theorem}{Theorem}
\newtheorem{axiom}{Axiom}
\newtheorem{proposition}{Proposition}
\newtheorem{lemma}{Lemma}

\title{\textbf{Fixed-point existence of self-consistent semiclassical gravity}\\[6pt]
{\large DRAFT}}

\author{Will Regelmann\thanks{Corresponding author.} \quad Claude (Anthropic)\thanks{AI
assistant; contributed to mathematical development and analysis.}}

\date{February 2026}

\begin{document}
\maketitle

\begin{abstract}
The semiclassical Einstein equation $G_{\mu\nu}[g] = 8\pi G\langle\hat{T}_{\mu\nu}\rangle_g$
is a fixed-point condition: the geometry and quantum state must be mutually consistent.
We address the question of whether such fixed points exist. We establish existence at three
levels: (1)~exactly, in the cosmological case via the trace anomaly
(Starobinsky, 1980); (2)~perturbatively near any classical solution, via a Banach
contraction with constant $\kappa \sim (m/M_P)^2 \ll 1$ for massive fields (the massless case, relevant to photons and gluons, remains open; see Section~\ref{sec:open}), established
established via a Hadamard decomposition of the order-reduced response kernel into a local
differential operator and a non-local smooth kernel; and (3)~conditionally in the general
case via the Schauder fixed-point theorem on globally hyperbolic spacetimes with compact
Cauchy surfaces, under six stated assumptions of which five are established and one remains
an open conjecture. The Planck scale emerges as the natural validity boundary in all three
analyses. An effective initial-value formulation, valid order-by-order in $\kappa$, is
derived from the global fixed-point condition.

\bigskip\noindent\textit{Note:} This is a companion paper to ``Gaussian temporal profile
of gravitational decoherence from the Einstein-Langevin equation.'' The decoherence
prediction of that paper does not depend on the existence results derived here; the present
paper addresses the independent question of whether the semiclassical framework is
internally self-consistent.
\end{abstract}

%======================================================================
\section{Introduction}
%======================================================================

The semiclassical Einstein equation:
\begin{equation}
G_{\mu\nu}[g] = 8\pi G\,\langle\Psi|\hat{T}_{\mu\nu}[g]|\Psi\rangle
\label{eq:sce}
\end{equation}
is not a dynamical equation in the usual sense. It is a fixed-point condition: we seek the
geometry $g_{\mu\nu}$ such that the quantum state it hosts, evolving on it, sources exactly
that geometry. The position that eq.~\eqref{eq:sce} may be fundamental (rather than an
approximation to quantized gravity) dates to M\o ller (1962) and Rosenfeld (1963) and has
been developed by Oppenheim~\cite{oppenheim}. Whether this position is internally
consistent depends on whether fixed points exist.

This paper establishes existence at three levels of generality, with explicit
quantitative control in the perturbative regime. The Planck scale appears in all three
analyses as the natural boundary of validity---not imposed by hand, but arising from the
mathematics of fixed-point convergence.

\medskip\noindent The perturbative contraction estimate requires a mass gap and therefore applies only to massive fields; extension to massless fields such as the photon and gluon is an open problem (Section~\ref{sec:open}).

%======================================================================
\section{The self-consistency map}
\label{sec:map}
%======================================================================

Define $\mathcal{F}: g_{\mu\nu} \mapsto g'_{\mu\nu}$ by:
\begin{enumerate}
\item Given trial geometry $g$, construct the QFT on $(\mathcal{M},g)$ and identify the
physical state $|\Psi_g\rangle$ satisfying the Hadamard condition.
\item Compute $S_{\mu\nu}[g] \equiv 8\pi G\langle\hat{T}_{\mu\nu}\rangle_{g,\text{ren}}$.
\item Solve $G_{\mu\nu}[g'] = S_{\mu\nu}[g]$ for $g'$.
\end{enumerate}
We seek $g^*$ such that $\mathcal{F}(g^*) = g^*$.

\medskip\noindent\textbf{State selection.} Step~1 requires care: the Hadamard condition
constrains the wavefront set of the two-point function but does not uniquely select a
state. In the cosmological case (Section~\ref{sec:starobinsky}), de~Sitter symmetry
selects the conformal vacuum. In the perturbative case (Section~\ref{sec:banach}), the
state on the corrected geometry is the unique perturbation of the state on the classical
background. For the Schauder argument (Section~\ref{sec:schauder}), multi-valuedness is
handled by restricting to a suitable branch or by applying Kakutani's theorem for
set-valued maps.

\medskip\noindent\textbf{Single-valuedness assumption.} The Banach argument (Section~\ref{sec:banach}) assumes $\mathcal{F}$ is single-valued: in the perturbative regime, the Hadamard state on the corrected geometry is the unique perturbation of the state on the classical background, so no state-selection ambiguity arises. The Schauder argument (Section~\ref{sec:schauder}) likewise assumes single-valuedness via the branch restriction noted above. An extension to the genuinely set-valued case via Kakutani's fixed-point theorem is possible but is not pursued here.

%======================================================================
\section{Cosmological fixed point: exact solution}
\label{sec:starobinsky}
%======================================================================

\begin{theorem}[Starobinsky, 1980 \cite{starobinsky}]
The self-consistency equation on FRW spacetime with conformal matter admits an exact
de~Sitter solution $a(t) = a_0 e^{H_0 t}$ with:
\begin{equation}
H_0^2 = \frac{180\pi}{G|a_2|}
\end{equation}
where $a_2$ is the coefficient of the Euler density in the trace anomaly. This solution is
a stable attractor: perturbations $\delta H$ are exponentially damped.
\end{theorem}

\noindent\emph{Proof sketch.} For constant $H$, the trace anomaly~\cite{capper_duff}
reduces to $\langle\hat{T}^\mu{}_\mu\rangle = -a_2 H_0^4/120\pi^2$. The Friedmann
equation yields $H_0$. Stability follows from the linearized equation for
$\delta H$~\cite{starobinsky}. \hfill$\square$

This constitutes an existence proof: the de~Sitter geometry, used to compute the quantum
stress-energy of fields living on it, sources exactly itself.

%======================================================================
\section{Perturbative fixed point: Banach contraction}
\label{sec:banach}
%======================================================================

Let $\bar{g}$ be a solution of the classical Einstein equations. Define the
quantum-corrected map:
\begin{equation}
\mathcal{F}(g) = \bar{g} + (G^{\bar{g}}_{\text{lin}})^{-1}\!\left[8\pi G\langle\hat{T}_{\mu\nu}\rangle^{(\text{quantum})}_g\right]
\end{equation}

\noindent We first establish a bound on the order-reduced response kernel, then use it
to prove the contraction estimate.

\begin{lemma}[Order-reduced response kernel bound]
\label{lem:kernel_bound}
Let $\phi$ be a massive scalar field with mass $m > 0$ on a globally hyperbolic spacetime
$(\mathcal{M},\bar{g})$ with compact Cauchy surface~$\Sigma$. After Parker--Simon order
reduction~\emph{\cite{parker_simon}}, the order-reduced response kernel
$\mathcal{K}^{\mathrm{red}}_{\mu\nu\alpha\beta}(x,x')
\equiv \delta\langle\hat{T}_{\mu\nu}\rangle^{\mathrm{red}}_{\mathrm{ren}}/\delta g^{\alpha\beta}(x')$
satisfies
\begin{equation}
\bigl\|\mathcal{K}^{\mathrm{red}}\bigr\|_{H^s \to H^{s-2}}
\;\leq\; C_{\mathcal{K}} \cdot \frac{\hbar\, m^2}{(4\pi)^2}
\label{eq:kernel_bound}
\end{equation}
uniformly for all $g \in K_\rho$, where $C_{\mathcal{K}}$ is a dimensionless constant
depending on $\rho$, $s$, and the background geometry~$\bar{g}$.
\end{lemma}

\noindent\emph{Proof.} We decompose the response kernel via the Hadamard
parametrix into a local differential operator and a non-local smooth kernel,
bound each separately, and combine via the triangle inequality.

\smallskip\noindent\textit{Step~1: Hadamard decomposition of the response kernel.}
By assumption~(A1), the renormalized stress-energy tensor is obtained by
point-splitting with Hadamard subtraction~\cite{hollands_wald_wick}:
\begin{equation}
\langle\hat{T}_{\mu\nu}(x)\rangle_{\mathrm{ren}}
= \lim_{x'\to x}\mathcal{D}_{\mu\nu}(x,x')\,W(x,x') + C_{\mu\nu}[g](x)
\label{eq:hadamard_ren}
\end{equation}
where $W = W_2 - H$ is the smooth Hadamard remainder ($W_2$ the full
two-point function, $H$ the universal Hadamard singular part determined by
local geometry via the Hadamard recursion), $\mathcal{D}_{\mu\nu}$ is the
second-order point-splitting differential operator, and $C_{\mu\nu}[g]$ collects
the finite local curvature counterterms (involving $R$,
$R_{\mu\alpha\nu\beta}R^{\alpha\beta}$, $\Box R$).

The functional derivative $\delta/\delta g^{\alpha\beta}(y)$ of
eq.~\eqref{eq:hadamard_ren} yields three types of contributions to the
response kernel $\mathcal{K}_{\mu\nu\alpha\beta}(x,y)$:
\begin{enumerate}
\item[(I)] Varying the smooth remainder~$W$ --- non-local, mediated by the
massive field propagator;
\item[(II)] Varying the point-splitting operator $\mathcal{D}_{\mu\nu}$ ---
local, from Christoffel symbol variation at $y=x$;
\item[(III)] Varying the curvature counterterms $C_{\mu\nu}$ --- local.
\end{enumerate}

\smallskip\noindent\textit{Step~2: Order reduction and local/non-local decomposition.}
Before order reduction, Term~(III) involves up to four derivatives of the
metric perturbation (from the $\Box R$ counterterm --- the source of the
Horowitz runaway instability~\cite{horowitz}). After Parker--Simon order
reduction~\cite{parker_simon} (assumption~(A5)), the fourth-derivative
terms are perturbatively eliminated using the classical Einstein equation,
leaving Term~(III) as a differential operator of order~$\leq 2$.

Term~(II) is a differential operator of order~$\leq 1$: varying
covariant derivatives introduces at most one derivative of~$\delta g$
via the Christoffel symbols. Term~(I) decomposes into contact
contributions at $y = x$ (absorbed into the local part) and a smooth
non-local kernel for $y\neq x$.

Define $\mathcal{K}^{\mathrm{red}}_{\mathrm{loc}}$ as the sum of all
local (distributional) contributions: a differential operator of
order~$\leq 2$ on~$\Sigma$. Define $\mathcal{K}_{\mathrm{nl}}$ as the
remaining non-local smooth kernel.

\smallskip\noindent\textit{Step~3: Bound on the local part.}
The local part takes the form
$\mathcal{K}^{\mathrm{red}}_{\mathrm{loc}}
= \sum_{|\gamma|\leq 2}a_\gamma(x)\,\partial^\gamma$
with smooth coefficients whose $L^\infty(\Sigma)$ norms are set by the
one-loop scale: $\|a_\gamma\|_{L^\infty} \leq C_{\mathrm{loc}}
\cdot\hbar m^2/(4\pi)^2$. A differential operator of order~$d$ with
$L^\infty$ coefficients maps $H^s(\Sigma)\to H^{s-d}(\Sigma)$
boundedly (by the Leibniz rule and Sobolev embedding on compact
manifolds). Therefore:
\begin{equation}
\bigl\|\mathcal{K}^{\mathrm{red}}_{\mathrm{loc}}\bigr\|_{H^s\to H^{s-2}}
\;\leq\; C_{\mathrm{loc}}\cdot\frac{\hbar\,m^2}{(4\pi)^2}.
\label{eq:loc_bound}
\end{equation}

\smallskip\noindent\textit{Step~4: Bound on the non-local part.}
For a massive field ($m > 0$), the smooth remainder~$W$ depends on the
metric through the Klein--Gordon equation
$(\Box_g + m^2 + \xi R)\phi = 0$. Its functional derivative
$\delta W/\delta g^{\alpha\beta}(y)$ is mediated by the massive
Klein--Gordon resolvent on~$\Sigma$, whose integral kernel decays
exponentially: $|(-\Delta_\Sigma + m^2 + \xi R|_\Sigma)^{-1}(x,y)|
\leq C\,e^{-m\cdot d_\Sigma(x,y)}$ for $d_\Sigma(x,y) \gg 1/m$. This
is a standard property of elliptic resolvents with positive mass on
compact Riemannian manifolds.

Consequently, $\mathcal{K}_{\mathrm{nl}}$ has a smooth kernel on
$\Sigma\times\Sigma$ with exponential decay. A smooth integral operator
on a compact manifold is a smoothing operator, bounded
$H^s\to H^{s'}$ for all~$s,s'$; its $H^s\to H^{s-2}$ operator norm
is bounded by its Hilbert--Schmidt norm:
\begin{equation}
\bigl\|\mathcal{K}_{\mathrm{nl}}\bigr\|_{H^s\to H^{s-2}}
\;\leq\; \|K_{\mathrm{nl}}\|_{L^2(\Sigma\times\Sigma)}
\;\leq\; C_{\mathrm{nl}}\cdot\frac{\hbar\,m^2}{(4\pi)^2}.
\label{eq:nl_bound}
\end{equation}

\smallskip\noindent\textit{Step~5: Combined bound.}
By the triangle inequality,
$\|\mathcal{K}^{\mathrm{red}}\|_{H^s\to H^{s-2}}
\leq \|\mathcal{K}^{\mathrm{red}}_{\mathrm{loc}}\|
+ \|\mathcal{K}_{\mathrm{nl}}\|
\leq C_{\mathcal{K}}\cdot\hbar m^2/(4\pi)^2$,
establishing~\eqref{eq:kernel_bound}.

\smallskip\noindent\textit{Step~6: Uniformity over $K_\rho$.}
By assumption~(A2), the map $g \mapsto \mathcal{K}^{\mathrm{red}}[g]$ is continuous.
The set $K_\rho$ is compact in $H^{s'}$ for $s' < s$ by Rellich--Kondrachov
(established in Section~\ref{sec:schauder}). A continuous function on a compact set
attains its supremum, so
$\sup_{g \in K_\rho}\|\mathcal{K}^{\mathrm{red}}[g]\|_{\mathrm{op}} < \infty$.
The constant $C_{\mathcal{K}}$ absorbs this supremum. \hfill$\square$

\smallskip\noindent\textit{Remark on the massless case.} The mass gap
$m > 0$ is essential in two places: it ensures exponential decay of the
non-local kernel~(Step~4), and it controls the infrared behavior of the
local coefficients~(Step~3). For massless or conformally coupled fields, the
non-local kernel decays only algebraically, the Hilbert--Schmidt bound may
diverge, and the argument breaks down. The massless case remains open.

\smallskip\noindent\textit{Remark on the K\"all\'en--Lehmann approach.}
On Minkowski spacetime, the bound can alternatively be established via the
K\"all\'en--Lehmann spectral representation of the stress-energy two-point
function~\cite{phillips_hu,hu_verdaguer}: the spectral gap above $4m^2$
provides polynomial suppression at high momenta, and the Schur test
translates this to an $H^s\to H^{s-2}$ bound. This approach is rigorous on
flat backgrounds but does not extend directly to general curved spacetimes,
where no K\"all\'en--Lehmann representation for the response kernel is
available. The Hadamard decomposition above avoids this difficulty.

\smallskip\noindent\textit{Remark on related work.} Meda, Pinamonti, and
Siemssen~\cite{meda_pinamonti_siemssen} established a Banach contraction for the
semiclassical Einstein equation on FLRW spacetimes, using symmetry-reduced bounds
specific to cosmological backgrounds. Gottschalk and Siemssen~\cite{gottschalk_siemssen}
showed that massive fields give polynomial decay of perturbations. The present Lemma
addresses the general (non-symmetric) case via the Hadamard decomposition.

\medskip

\begin{theorem}[Perturbative existence---massive fields]
\label{thm:banach}
Let $\bar{g}$ be a classical solution with sub-Planckian curvature
$R \ll \ell_P^{-2}$, and let the matter content consist of massive fields with
$m > 0$. Under the hypotheses of Lemma~\ref{lem:kernel_bound}, the self-consistency
map $\mathcal{F}$ is a contraction in a neighborhood of $\bar{g}$ in $H^s$, with
Lipschitz constant
\begin{equation}
\kappa \;=\; 8\pi G \cdot C_{\mathrm{inv}} \cdot
\sup_{g \in K_\rho}\bigl\|\mathcal{K}^{\mathrm{red}}[g]\bigr\|_{H^s \to H^{s-2}}
\;\leq\; C \cdot \frac{1}{2\pi}\!\left(\frac{m}{M_P}\right)^{\!2}\ell_P^2\, R
\;\ll\; 1
\label{eq:kappa}
\end{equation}
where $C_{\mathrm{inv}}$ is the norm of the inverse linearized Einstein operator
from~\emph{(A4)} and $C$ is a dimensionless constant. By the Banach fixed-point
theorem~\cite{banach}, there exists a unique self-consistent solution $g^*$ in a neighborhood of
$\bar{g}$, and the iteration $g_{n+1} = \mathcal{F}(g_n)$ converges exponentially:
$\|g_n - g^*\|_{H^s} \leq \kappa^n\|g_0 - g^*\|_{H^s}$.
\end{theorem}

\noindent\emph{Proof.} The Lipschitz constant of $\mathcal{F}$ is
\begin{equation}
\kappa = 8\pi G\,\bigl\|(G_{\mathrm{lin}})^{-1}\bigr\|_{H^{s-2}\to H^s}
\cdot \sup_{g \in K_\rho}
\bigl\|\mathcal{K}^{\mathrm{red}}[g]\bigr\|_{H^s \to H^{s-2}}.
\end{equation}
By assumption~(A4), the inverse linearized Einstein operator satisfies
$\|(G_{\mathrm{lin}})^{-1}\|_{H^{s-2}\to H^s} \leq C_{\mathrm{inv}}$, with
$C_{\mathrm{inv}} \sim 1/R$, since the harmonic gauge condition~(A4) renders the linearized Einstein operator elliptic with principal symbol determined by the background curvature scale~$R$; standard elliptic regularity estimates then give an inverse operator norm scaling as~$1/R$. By Lemma~\ref{lem:kernel_bound}, the response
kernel norm is bounded by $C_{\mathcal{K}} \cdot \hbar m^2/(4\pi)^2$. Combining:
\begin{equation}
\kappa \;\leq\; 8\pi G \cdot \frac{C_{\mathrm{inv}}C_{\mathcal{K}}\,\hbar\, m^2}{(4\pi)^2}
\;\sim\; \frac{\hbar G\, m^2}{2\pi}
\cdot \frac{1}{R}
\;=\; \frac{1}{2\pi}\!\left(\frac{m}{M_P}\right)^{\!2}\ell_P^2\, R.
\end{equation}
For sub-Planckian curvature ($R \ll \ell_P^{-2}$, i.e., $\ell_P^2 R \ll 1$)
and any known particle mass ($m \ll M_P$), we have $\kappa \ll 1$.
The Banach fixed-point theorem~\cite{banach} then guarantees a unique
fixed point in the closed ball around $\bar{g}$ of radius
$\rho/(1-\kappa)$, with geometric convergence.

\hfill$\square$

\medskip

For all known particles, $\kappa_{\text{mass}} \equiv (m/M_P)^2 \ll 1$:
\begin{center}
\begin{tabular}{lcc}
\toprule
Particle & Mass (kg) & $\kappa_{\text{mass}} = (m/M_P)^2$ \\
\midrule
Electron & $9.1\times10^{-31}$ & $\sim 10^{-45}$ \\
Proton & $1.7\times10^{-27}$ & $\sim 10^{-39}$ \\
Top quark & $3.1\times10^{-25}$ & $\sim 10^{-34}$ \\
\bottomrule
\end{tabular}
\end{center}

%======================================================================
\section{Non-perturbative existence via Schauder}
\label{sec:schauder}
%======================================================================

\subsection{Setup}

Let $(\mathcal{M},\bar{g})$ be a globally hyperbolic spacetime with \emph{compact} Cauchy
surface $\Sigma$. Fix $s > n/2 + 2 = 4$ and define the convex compact set:
\begin{equation}
K_\rho \equiv \bigl\{h \in H^s(\Sigma) : \|h-\bar{h}\|_{H^s} \leq \rho\bigr\}
\cap \mathcal{G}_{\text{harm}} \cap \mathcal{L}
\label{eq:K_rho}
\end{equation}
where $\mathcal{G}_{\text{harm}}$ is the affine set of metrics in linearized harmonic
gauge and $\mathcal{L}$ is the open set of Lorentzian-signature metrics. By the
Rellich--Kondrachov theorem on compact $\Sigma$, $K_\rho$ is compact in $H^{s'}$ for
any $s' < s$.

\subsection{Assumptions}

\begin{itemize}
\item[\textbf{(A1)}] \textit{Hadamard state existence.} For each $g\in K_\rho$, a
Hadamard state $|\Psi_g\rangle$ exists and $\langle\hat{T}_{\mu\nu}\rangle_{g,\text{ren}}$
is well-defined. \textbf{[Established:} Radzikowski~\cite{radzikowski},
Fulling--Sweeny--Wald~\cite{fulling_sweeny_wald}.\textbf{]}

\item[\textbf{(A2)}] \textit{Continuous metric dependence.} The map $g\mapsto
\langle\hat{T}_{\mu\nu}\rangle_{g,\text{ren}}$ is continuous $H^s\to H^{s-2}$.
\textbf{[Established:} Hollands--Wald~\cite{hollands_wald_wick}.\textbf{]}

\item[\textbf{(A3)}] \textit{Uniform stress-energy bound.} There exists $B(\rho)<\infty$
with $\|\langle\hat{T}_{\mu\nu}\rangle_{g,\text{ren}}\|_{H^{s-2}} \leq B(\rho)$ for all
$g\in K_\rho$. \textbf{[Conjecture.} Established for FRW spacetimes
\cite{pinamonti_siemssen}; open in general.\textbf{]}

\item[\textbf{(A4)}] \textit{Linearized Einstein well-posedness.} The linearized Einstein
equation in harmonic gauge has a unique solution $h\in H^s$ with
$\|h\|_{H^s} \leq C_{\text{inv}}\|S\|_{H^{s-2}}$.
\textbf{[Established:} Choquet-Bruhat~\cite{choquet_bruhat}.\textbf{]}

\item[\textbf{(A5)}] \textit{Order reduction.} Fourth-derivative terms from the trace
anomaly can be perturbatively eliminated. \textbf{[Established:} Parker--Simon
\cite{parker_simon}.\textbf{]}

\item[\textbf{(A6)}] \textit{Uniform contraction bound.}
(Included for completeness; used only in the Banach uniqueness argument of Section~\ref{sec:banach}, not in the Schauder existence argument below.)\newline
$\kappa_{\sup} \equiv 8\pi G\,C_{\text{inv}}\cdot\sup_{g\in K_\rho}\|\delta\langle\hat{T}_{\mu\nu}\rangle_{\text{ren}}/\delta g\|_{\text{op}} < 1$.
\textbf{[Established for massive fields.} Lemma~\ref{lem:kernel_bound} establishes the
operator-norm bound with constant $\kappa \sim (m/M_P)^2\ell_P^2 R \ll 1$ for
sub-Planckian curvature. Open for massless fields.\textbf{]}
\end{itemize}

\subsection{Theorem}

\begin{theorem}[Non-perturbative existence (conditional on (A3))]
\label{thm:schauder}
Let $(\mathcal{M},\bar{g})$ be globally hyperbolic with compact Cauchy surface $\Sigma$,
and $\bar{g}$ a classical solution in linearized harmonic gauge. Under assumptions
\emph{(A1)--(A6)}, the self-consistency map $\mathcal{F}$ has at least one fixed point
$g^* \in K_\rho$ in the $H^{s'}$ topology for any $s' < s$.
\end{theorem}

\noindent\emph{Proof.} We verify the hypotheses of the Schauder fixed-point theorem
\cite{schauder}.

\smallskip\noindent\textit{Step 1: $K_\rho$ is convex and compact.} $K_\rho$ is the
intersection of a closed $H^s$ ball with the convex affine sets $\mathcal{G}_{\text{harm}}$
and $\mathcal{L}$, hence convex and closed in $H^s$. By Rellich--Kondrachov, it is compact
in $H^{s'}$.

\smallskip\noindent\textit{Step 2: $\mathcal{F}$ maps $K_\rho$ into $K_\rho$.} Given
$g\in K_\rho$: (A1) provides a Hadamard state with well-defined renormalized
stress-energy; (A5) reduces the fourth-order system to second order; (A4) gives
$\mathcal{F}(g) \in H^s$ with $\|\mathcal{F}(g)-\bar{h}\|_{H^s} \leq
C_{\text{inv}}\|S_{\mu\nu}\|_{H^{s-2}} \leq C_{\text{inv}}B(\rho)$ by (A3). Choose
$\rho = C_{\text{inv}}B(\rho)$ (a self-consistency condition on $\rho$ that is satisfiable
since $C_{\text{inv}}B(\rho)\to 0$ as $\rho\to 0$ for perturbative regimes). The harmonic
gauge and Lorentzian conditions are preserved by the linearized evolution.

\smallskip\noindent\textit{Step 3: $\mathcal{F}$ is continuous.} By (A2), $g\mapsto
S_{\mu\nu}[g]$ is continuous in $H^{s-2}$; by (A4), the Einstein solve is continuous;
composing, $\mathcal{F}: K_\rho\to K_\rho$ is continuous in $H^{s'}$.

\smallskip\noindent By the Schauder theorem~\cite{schauder}, $\mathcal{F}$ has a fixed
point in $K_\rho$.\hfill$\square$

\noindent\emph{Remark.} Assumption~(A6) is not used in the Schauder argument (which
requires only continuity, not contractivity). It would be needed for a Banach fixed-point
argument guaranteeing uniqueness and constructive convergence.

%======================================================================
\section{Effective initial-value formulation}
\label{sec:ivp}
%======================================================================

Theorem~\ref{thm:schauder} establishes global existence. Practical physics requires a
Cauchy formulation: given data on a spacelike hypersurface, what is the evolution? The
global fixed-point condition does not directly reduce to a Cauchy problem because the
Hadamard state is defined by a global microlocal condition. We derive an effective
formulation valid order-by-order in $\kappa$.

At order $\kappa^0$, the evolution is classical GR with given matter initial data.

At order $\kappa^1$, the semiclassical correction $\langle\hat{T}_{\mu\nu}\rangle^{(1)}$
is computed on the order-$\kappa^0$ background geometry. By Parker--Simon~\cite{parker_simon}
order reduction, the resulting system is second-order in time, with initial data for
$g_{\mu\nu}$ and $\partial_t g_{\mu\nu}$ on $\Sigma$ supplemented by the adiabatic vacuum
state at each order.

The order-reduced system at each order $\kappa^n$ is a second-order hyperbolic PDE,
well-posed by Choquet-Bruhat~\cite{choquet_bruhat}. The perturbative expansion is valid
while $\kappa \ll 1$, i.e., throughout the sub-Planck regime.

\noindent\textbf{Open problem.} Non-perturbative well-posedness of the order-reduced
semiclassical Einstein equation on general globally hyperbolic spacetimes without symmetry
assumptions has been conjectured by Ju\'arez-Aubry et al.~\cite{juarez_aubry} and
established for cosmological spacetimes~\cite{pinamonti_siemssen}. The general case
remains open.

%======================================================================
\section{Planck-scale validity boundary}
\label{sec:planck}
%======================================================================

The contraction constant~\eqref{eq:kappa} satisfies $\kappa \to 1$ when
$R \sim \ell_P^{-2}$, i.e., at Planck-scale curvatures. This identifies the natural
validity boundary of the framework: not imposed by hand, but derived from the requirement
$\kappa < 1$ for the fixed-point iteration to converge.

The framework also fails for observables that require summing over spacetime topologies
(e.g., the Page curve of black hole evaporation). Even when $\kappa \ll 1$---so the
perturbative expansion is valid---certain correlators receive $\mathcal{O}(1)$
contributions from topology-changing saddles absent from the fixed-topology
description. This is a distinct limitation from the Planck-scale breakdown above, and is
not self-diagnosable within the framework.

%======================================================================
\section{Open problems}
\label{sec:open}
%======================================================================

\begin{enumerate}

\item \textbf{Extension to massless fields.}
Lemma~\ref{lem:kernel_bound} establishes the $H^s \to H^{s-2}$ bound for massive fields
via the Hadamard decomposition: the mass gap $m > 0$ ensures exponential decay of the
non-local kernel and controls the infrared behavior of the local coefficients. Extending
the bound to massless or conformally coupled fields requires new IR control, since the
non-local kernel decays only algebraically when $m = 0$ and the Hilbert--Schmidt bound
may diverge.

\item \textbf{Uniform stress-energy bound (A3).} The Pinamonti--Siemssen
result~\cite{pinamonti_siemssen} establishes this for FRW spacetimes. Extension to general
globally hyperbolic spacetimes with compact Cauchy surfaces requires controlling the
state-dependent part of $\langle\hat{T}_{\mu\nu}\rangle_{\text{ren}}$ uniformly over
$K_\rho$---an open problem.

\item \textbf{Relaxing the compact Cauchy surface assumption.} The Schauder argument
requires compactness of $K_\rho$, ensured by Rellich--Kondrachov on compact $\Sigma$.
Extension to asymptotically flat spacetimes requires weighted Sobolev spaces and a
different compactness argument.

\item \textbf{Non-perturbative IVP well-posedness.} See Section~\ref{sec:ivp}.

\end{enumerate}

%======================================================================
\begin{thebibliography}{99}

\bibitem{oppenheim} J.~Oppenheim, A postquantum theory of classical gravity?,
\textit{Phys.\ Rev.\ X}\ \textbf{13}, 041040 (2023).

\bibitem{starobinsky} A.~A.~Starobinsky, A new type of isotropic cosmological models
without singularity, \textit{Phys.\ Lett.\ B}\ \textbf{91}, 99 (1980).

\bibitem{capper_duff} D.~M.~Capper and M.~J.~Duff, Trace anomalies in dimensional
regularization, \textit{Nuovo Cimento A}\ \textbf{23}, 173 (1974).

\bibitem{horowitz} G.~T.~Horowitz, Semiclassical relativity: The weak-field limit,
\textit{Phys.\ Rev.\ D}\ \textbf{21}, 1445 (1980).

\bibitem{phillips_hu} N.~G.~Phillips and B.~L.~Hu, Noise kernel in stochastic gravity and
stress energy bitensor, \textit{Phys.\ Rev.\ D}\ \textbf{63}, 104001 (2001).

\bibitem{hu_verdaguer} B.~L.~Hu and E.~Verdaguer, Stochastic gravity: Theory and
applications, \textit{Living Rev.\ Relativ.}\ \textbf{11}, 3 (2008).

\bibitem{radzikowski} M.~J.~Radzikowski, Micro-local approach to the Hadamard condition,
\textit{Commun.\ Math.\ Phys.}\ \textbf{179}, 529 (1996).

\bibitem{fulling_sweeny_wald} S.~A.~Fulling, M.~Sweeny, and R.~M.~Wald, Singularity
structure of the two-point function in quantum field theory in curved spacetime,
\textit{Commun.\ Math.\ Phys.}\ \textbf{63}, 257 (1978).

\bibitem{hollands_wald_wick} S.~Hollands and R.~M.~Wald, Local Wick polynomials and time
ordered products of quantum fields in curved spacetime, \textit{Commun.\ Math.\ Phys.}\
\textbf{223}, 289 (2001).

\bibitem{choquet_bruhat} Y.~Choquet-Bruhat, Th\'eor\`eme d'existence pour certains
syst\`emes d'\'equations aux d\'eriv\'ees partielles non lin\'eaires, \textit{Acta Math.}\
\textbf{88}, 141 (1952).

\bibitem{banach} S.~Banach, Sur les op\'erations dans les ensembles abstraits et leur
application aux \'equations int\'egrales, \textit{Fund.\ Math.}\ \textbf{3}, 133 (1922).

\bibitem{schauder} J.~Schauder, Der Fixpunktsatz in Funktionalr\"aumen,
\textit{Studia Math.}\ \textbf{2}, 171 (1930).

\bibitem{parker_simon} L.~Parker and J.~Z.~Simon, Einstein equation with quantum
corrections reduced to second order, \textit{Phys.\ Rev.\ D}\ \textbf{47}, 1339 (1993).

\bibitem{pinamonti_siemssen} N.~Pinamonti and D.~Siemssen, Global existence of solutions
of the semiclassical Einstein equation for cosmological spacetimes,
\textit{Commun.\ Math.\ Phys.}\ \textbf{334}, 171 (2015).

\bibitem{juarez_aubry} B.~A.~Ju\'arez-Aubry, B.~S.~Kay, T.~Miramontes, and D.~Sudarsky,
On the initial value problem for semiclassical gravity without and with quantum state
collapses, \textit{J.\ Cosmol.\ Astropart.\ Phys.}\ \textbf{2023}(01), 040 (2023).

\bibitem{meda_pinamonti_siemssen} P.~Meda, N.~Pinamonti, and D.~Siemssen,
Existence and uniqueness of solutions of the semiclassical Einstein equation in
cosmological models, \textit{Ann.\ Henri Poincar\'e}\ \textbf{22}, 3965 (2021).

\bibitem{gottschalk_siemssen} H.~Gottschalk and D.~Siemssen,
The cosmological semiclassical Einstein equation: Stability of the de~Sitter
spacetime, preprint arXiv:2201.10288 (2022).

\end{thebibliography}

\end{document}
